\subsubsection*{命题 6.2.}
设 $X$ 为代数闭域 $k$ 上的 $n$ 维光滑连通射影簇. 设 $D$ 为 Lefschetz 线束, $\eta$ 为其泛点, $\bar\eta$ 为 $\eta$ 上的一个几何点. 则线束 $D$ 满足条件 (A) 当且仅当 $D\cap\check X$ 中每个点处的惯性群在 $H^{n-1}(X_{\bar\eta},\mathbb Q_\ell)$ 上的作用非平凡.

\subsubsection*{定理 6.3.}
在 6.2 的假设下, 下列断言成立.

\medskip
\noindent
\textup{(6.3.1)} 若 $n=\dim X$ 为奇数, 则 $X$ 的每个 Lefschetz 线束都满足条件 (A).

\medskip
\noindent
\textup{(6.3.2)} 假设 $n=\dim X$ 为偶数, 或 $\operatorname{char}(k)\ne2$. 则下列条件等价:
\begin{enumerate}
\item[(i)] $D$ 不满足 (5.3.5) 中的条件 (A);
\item[(ii)] $\pi_1(D-D\cap\check X)$ 在 $H^{n-1}(X_{\bar\eta},\mathbb Q_\ell)$ 上的作用平凡, 且
  $D\cap\check X\ne\varnothing$;
\item[(iii)] $D\cap\check X$ 中存在一点, 其惯性群在 $H^{n-1}(X_{\bar\eta},\mathbb Q_\ell)$ 上的作用平凡;
\item[(iv)] 层 $R^{n-1}\rho_*\mathbb Q_\ell$ 在 $D$ 上为常层, 且 $D\cap\check X$ 非空;
\item[(v)] 限制 $R^{n-1}\rho_*\mathbb Q_\ell\vert_U$ 在 $U$ 上为常层, 其中 $U$ 为 $D$ 的一个非空开子集, $\rho$ 在其上光滑, 且 $\rho$ 并非光滑态射.
\end{enumerate}

\noindent
\textup{(6.3.3)} 若再假设 $X$ 满足 (5.2.2) 中的条件 (LV), 则条件 (i)--(v) 也等价于:
\begin{enumerate}
\item[(vi)] 限制映射
\[
H^{n-1}(X,\mathbb Q_\ell)
\longrightarrow H^{n-1}(X_{\bar\eta},\mathbb Q_\ell)
\]
为满射 (因而由弱 Lefschetz 定理为同构).
\end{enumerate}

\subsubsection*{证明.}
断言 6.3.1 由 XV 3.4 得出. 关于 6.3.2 中的等价性, 由 6.2 得 (i) 与 (iii) 等价, 由 6.1.10 得 (ii) 与 (iii) 等价, (ii) 与 (v) 按定义即等价, 而由 5.3.4 得 (iv) 与 (v) 等价. 最后, 若 $X$ 满足条件 (LV), 则 5.6.10 给出 (ii) 与 (vi) 的等价性.

\subsubsection*{注 6.3.4.}
条件 (A) 不依赖于 Lefschetz 线束 $D$ 的选取, 而只依赖于所选的 (Lefschetz) 浸入 $i:X\longrightarrow\mathbb P^r$. 当 $\dim X$ 为奇数时, 这由 (6.3.1) 显然可见. 当 $\dim X$ 为偶数且 $D\cap\check X\ne\varnothing$ 时, 不满足条件 (A) 等价于由群 $H^{n-1}(X\cap H_s,\mathbb Q_\ell)$ 构成的 $\check{\mathbb P}^{r}-\check X$ 上的局部系统为常系统. 这由判别准则 (ii)、所论局部系统在 $\check X$ 的泛点处驯分歧这一事实, 以及 (6.1.1) 中的满射性得出. 若 $D\cap\check X=\varnothing$, 则条件 (A) 自动成立.

\begin{editorialnote}
原印本及英译底本在 6.3 的证明中将 (ii)$\Longleftrightarrow$(iii) 指向 (6.1.6); 这里应使用 6.1.10. 注 6.3.4 又将条件 (A) 本身说成局部系统为常系统, 与紧前的判别准则 (6.3.2)(i)、(ii) 相反: 当奇异成员非空时, 应为\emph{不满足} (A) 等价于常系统. 当奇异成员为空时, (A) 为空条件, 自动成立, 不能沿用该等价说法. 此处明确区分这两种情形, 并将所需满射的引用由 (6.1.2) 改为 (6.1.1); 参数空间写成对偶射影空间. 关于不依赖线束的结论不变: 对偶簇若非超曲面, Lefschetz 线束没有奇异成员; 若为超曲面, 则使用上述判别准则及满射.
\end{editorialnote}

\subsubsection*{推论 6.4.}
在 6.2 的假设下, 假设 $n=\dim X$ 为偶数或 $\operatorname{char}(k)\ne2$, 并且 $X$ 满足条件 (LV).\footnote{印本中连词的位置使条件 (LV) 仅附属于第二种选择. 证明引用了 (6.3.3) 中的等价关系 (i)$\Longleftrightarrow$(vi), 因而两种情形都需要条件 (LV).}
则存在整数 $d_o$, 使得当 $d\ge d_o$ 时, 由 $d$ 次超曲面在 $X$ 上截得的截面所成的每个 Lefschetz 线束都满足条件 (A).

\subsubsection*{证明.}
由 (6.3.3) 中的等价关系 (i)$\Longleftrightarrow$(vi), 只须证明: 若 $Y(d)$ 表示 $X$ 的一个 $d$ 次光滑超曲面截面, 则
\begin{equation}\tag{6.4.1}
\dim H^{n-1}(Y(d),\mathbb Q_\ell)\longrightarrow\infty
\qquad\text{当 }d\longrightarrow\infty.
\end{equation}
事实上, 有以下精确结果.

\begin{editorialnote}
原印本及英译底本在推论 6.4 的证明中引用 (6.3.2) 的 (vi), 但 (vi) 实际列于 (6.3.3), 此处订正. 引理 6.4.2 则承接证明中已选取的\emph{光滑}截面 $Y(d)$; 其独立陈述漏写了光滑性, 本文补明. 证明使用的伴随公式、Poincar\'e 对偶性及弱 Lefschetz 比较均针对该光滑截面. 若允许奇异截面, 相同次数的光滑三次平面曲线与结点三次平面曲线的一阶 Betti 数分别为二与一, 故原陈述不能成立.
\end{editorialnote}

\subsubsection*{引理 6.4.2.}
设 $X$ 为代数闭域 $k$ 上的 $n\ge1$ 维光滑连通射影簇, $Y(d)$ 表示一个 $d$ 次光滑超曲面截面. 则存在多项式 $f(T)\in\mathbb Z[T]$, 使得
\begin{align}
f(T)&=\deg(X)T^n+\text{次数至多为下列数的多项式:}n-1,
  \tag{6.4.2.1}\\
f(d)&=\dim H^{n-1}(Y(d),\mathbb Q_\ell)\qquad(d\ge1).
  \tag{6.4.2.2}
\end{align}

\subsubsection*{证明.}
由 XVII 5.7.5, 若 $C(X)$ 表示 $X$ 的总 Chern 类, $H$ 表示超平面的第一 Chern 类, 则 $Y(d)$ 的 Euler--Poincaré 示性数为
\begin{equation}\tag{6.4.2.3}
\chi(Y(d))
=\int_X\frac{dH\,C(X)}{1+dH}
=\sum_{i=1}^{n}(-1)^{i-1}d^i\int_X H^iC_{n-i}.
\end{equation}
另一方面, Poincaré 对偶性及弱 Lefschetz 定理给出
\begin{equation}\tag{6.4.2.4}
\chi(Y(d))
=(-1)^{n-1}\dim H^{n-1}(Y(d),\mathbb Q_\ell)
+2\sum_{i\ge1}(-1)^{n-1-i}
  \dim H^{n-1-i}(X,\mathbb Q_\ell).
\end{equation}
结合 (6.4.2.3) 和 (6.4.2.4) 即得结论.

